\documentclass{article}
\usepackage{amsmath,amssymb,amsthm}
\usepackage{algorithm,algorithmic}
\usepackage{hyperref}
\usepackage{enumitem}
\newtheorem{theorem}{Theorem}
\newtheorem{lemma}{Lemma}
\newtheorem{assumption}{Assumption}
\newtheorem{definition}{Definition}
\newcommand{\grad}{\nabla}
\newcommand{\E}{\mathbb{E}}
\newcommand{\R}{\mathbb{R}}

\begin{document}

\section*{Algorithm Name}
\textbf{Nesterov’s Accelerated Gradient with Polynomial Momentum Schedule}  
(the momentum coefficient is chosen as $\displaystyle \beta_k=\frac{k-1}{k+2}$).

\section*{Mathematical Formulation}
Let $f:\R^d\rightarrow\R$ be the objective function.  
Initialize $x_{0}=y_{0}\in\R^{d}$ and a stepsize $\alpha>0$.  
For $k=0,1,2,\dots$ perform
\begin{align}
    \beta_k &:= \frac{k-1}{k+2},\qquad\text{(momentum schedule)}\label{eq:beta}\\[4pt]
    y_k    &= x_k + \beta_k\,(x_k-x_{k-1}),\qquad\text{(extrapolation)}\label{eq:y}\\[4pt]
    x_{k+1}&= y_k - \alpha\,\grad f(y_k).\qquad\text{(gradient step)}\label{eq:x}
\end{align}
When $k=0$ we set $x_{-1}=x_0$ so that $\beta_0= -\frac{1}{2}$ yields $y_0=x_0$.

\section*{Pseudocode}
\begin{algorithm}[H]
\caption{Accelerated Gradient with $\beta_k=(k-1)/(k+2)$}
\begin{algorithmic}[1]
\STATE \textbf{Input:} stepsize $\alpha>0$, initial point $x_0\in\R^d$, maximum iterations $K$.
\STATE Set $x_{-1}\gets x_0$.
\FOR{$k=0$ \TO $K-1$}
    \STATE $\beta_k \gets (k-1)/(k+2)$
    \STATE $y_k \gets x_k + \beta_k (x_k - x_{k-1})$
    \STATE $g_k \gets \grad f(y_k)$
    \STATE $x_{k+1} \gets y_k - \alpha g_k$
\ENDFOR
\STATE \textbf{Output:} $x_K$ (or the best iterate in terms of $f$).
\end{algorithmic}
\end{algorithm}

\section*{Dry Run Example}
Consider the one–dimensional quadratic
\[
f(w) = (w-3)^2, \qquad \grad f(w)=2(w-3).
\]
Choose $\alpha=0.1$, $x_0=0$, and $K=3$.

\begin{tabular}{c|c|c|c|c}
$k$ & $\beta_k$ & $y_k$ & $g_k=\grad f(y_k)$ & $x_{k+1}=y_k-\alpha g_k$ \\ \hline
0 & $-1/2$ & $y_0 = 0$ & $g_0 = 2(0-3)=-6$ & $x_1 = 0 -0.1(-6)=0.6$ \\[4pt]
1 & $0$   & $y_1 = x_1$ & $g_1 = 2(0.6-3)=-4.8$ & $x_2 = 0.6 -0.1(-4.8)=1.08$ \\[4pt]
2 & $1/4$ & $y_2 = 1.08 + \tfrac14(1.08-0.6)=1.08+0.12=1.20$ &
$g_2 = 2(1.20-3)=-3.6$ & $x_3 = 1.20-0.1(-3.6)=1.56$
\end{tabular}

Objective values:
\[
f(x_0)=9,\quad f(x_1)=5.76,\quad f(x_2)=3.6864,\quad f(x_3)=2.0736.
\]
The sequence clearly descends toward the optimum $w^\star=3$.

\section*{Assumptions}
\begin{assumption}[Convexity]\label{ass:convex}
$f$ is convex, i.e. for all $u,v\in\R^d$,
\[
f(v)\ge f(u)+\langle \grad f(u),\,v-u\rangle .
\]
\end{assumption}
\begin{assumption}[Smoothness]\label{ass:smooth}
$f$ is $L$‑smooth: for all $u,v\in\R^d$,
\[
\|\grad f(u)-\grad f(v)\|\le L\|u-v\|
\quad\Longleftrightarrow\quad
f(v)\le f(u)+\langle\grad f(u),v-u\rangle+\frac{L}{2}\|v-u\|^2 .
\]
\end{assumption}
\begin{assumption}[Stepsize]\label{ass:step}
The stepsize satisfies $0<\alpha\le 1/L$.
\end{assumption}
All subsequent results hold under Assumptions~\ref{ass:convex}--\ref{ass:step}.

\section*{Main Convergence Theorem}
\begin{theorem}\label{thm:rate}
Let $\{x_k\}$ be generated by \eqref{eq:beta}--\eqref{eq:x} with $\alpha\le 1/L$.
Define the optimal value $f^\star:=\min_{w}f(w)$ and let $w^\star$ be any minimizer.
Then for every $k\ge 1$,
\[
f(x_k)-f^\star \;\le\; \frac{2L\|x_0-w^\star\|^2}{(k+1)^2}.
\tag{T1}
\]
Consequently $f(x_k)\to f^\star$ with the optimal accelerated rate $O(1/k^2)$.
\end{theorem}

\section*{Theoretical Properties}
\begin{itemize}[leftmargin=*]
\item \textbf{Stability.} The choice $\beta_k=(k-1)/(k+2)$ guarantees that $\beta_k\in[0,1)$ for $k\ge1$ and the extrapolation does not overshoot the previous iterate.
\item \textbf{Hyper‑parameter sensitivity.} The only free scalar is the stepsize $\alpha$. Under $\alpha\le 1/L$ the bound \eqref{T1} holds; larger $\alpha$ may destroy the $O(1/k^2)$ rate.
\item \textbf{Comparison to baselines.}
    \begin{enumerate}
        \item Gradient descent with stepsize $1/L$ yields $f(x_k)-f^\star\le \frac{L\|x_0-w^\star\|^2}{2k}$.
        \item The present method improves the $1/k$ rate to $1/k^2$ without any line‑search.
    \end{enumerate}
\item \textbf{Extension.} The same proof works for any $\beta_k$ generated by the recurrence
        $\beta_k=\frac{t_{k-1}-1}{t_k}$ with $t_k$ satisfying $t_{k+1}^2-t_{k+1}=t_k^2$, which yields the classic Nesterov schedule $t_k=\frac{k+2}{2}$.
\end{itemize}

\section*{Discussion}
The theorem shows that the polynomial momentum schedule reproduces the optimal accelerated convergence for convex smooth objectives. The proof relies solely on convexity, smoothness, and the algebraic identity linking $\beta_k$ and the auxiliary sequence $t_k=k+2$, thus the result is robust to problem dimension and does not require stochastic assumptions.

\section*{Preliminaries (Definitions and Lemmas)}
\begin{definition}[Auxiliary sequence]\label{def:tk}
Let $t_k:=k+2$ for $k\ge0$. Observe that
\[
\beta_k=\frac{k-1}{k+2}= \frac{t_{k-1}-1}{t_k}.
\tag{D1}
\]
\end{definition}
\begin{lemma}[Smoothness inequality]\label{lem:smooth}
For $L$‑smooth $f$ and any $u,v\in\R^d$,
\[
f(v)\le f(u)+\langle\grad f(u),v-u\rangle+\frac{L}{2}\|v-u\|^2 .
\tag{L1}
\]
\end{lemma}
\begin{lemma}[Three‑point identity]\label{lem:three}
For any $a,b,c\in\R^d$ and any $\lambda>0$,
\[
\|a-c\|^2 = \|a-b\|^2 + 2\langle a-b,c-b\rangle + \|c-b\|^2 .
\tag{L2}
\]
\end{lemma}

\section*{Main Proof (Step‑by‑Step Derivation)}
We prove Theorem~\ref{thm:rate}. Throughout the proof $w^\star$ denotes a minimizer of $f$.

\noindent\textbf{Step 1 (Potential function).}  
Define for $k\ge0$
\[
\Phi_k := t_k^2\bigl(f(x_k)-f^\star\bigr)+\frac{L}{2}\|x_k-w^\star\|^2 .
\tag{P1}
\]
Our goal is to show $\Phi_{k+1}\le \Phi_k$ for all $k\ge0$; telescoping then yields the bound.

\noindent\textbf{Step 2 (Relation between $x_{k+1}$ and $y_k$).}  
From \eqref{eq:x} we have
\[
x_{k+1}=y_k-\alpha\grad f(y_k).
\tag{S2}
\]

\noindent\textbf{Step 3 (Apply smoothness at $y_k$).}  
Using Lemma~\ref{lem:smooth} with $u=y_k$, $v=x_{k+1}$ and $\alpha\le 1/L$,
\[
f(x_{k+1})\le f(y_k)+\langle\grad f(y_k),x_{k+1}-y_k\rangle+\frac{L}{2}\|x_{k+1}-y_k\|^2 .
\]
Since $x_{k+1}-y_k=-\alpha\grad f(y_k)$,
\[
f(x_{k+1})\le f(y_k)-\alpha\|\grad f(y_k)\|^2+\frac{L\alpha^{2}}{2}\|\grad f(y_k)\|^2 .
\tag{S3}
\]
Because $\alpha\le 1/L$, $L\alpha^{2}/2\le \alpha/2$, thus
\[
f(x_{k+1})\le f(y_k)-\frac{\alpha}{2}\|\grad f(y_k)\|^2 .
\tag{S3'}
\]

\noindent\textbf{Step 4 (Convexity at $w^\star$).}  
Convexity \eqref{ass:convex} with $u=y_k$, $v=w^\star$ gives
\[
f^\star\ge f(y_k)+\langle\grad f(y_k),w^\star-y_k\rangle .
\tag{S4}
\]
Rearranging,
\[
\langle\grad f(y_k),y_k-w^\star\rangle\ge f(y_k)-f^\star .
\tag{S4'}
\]

\noindent\textbf{Step 5 (Combine \eqref{S3'} and \eqref{S4'}).}  
Multiply \eqref{S4'} by $\alpha$ and add to \eqref{S3'}:
\[
\begin{aligned}
f(x_{k+1})-f^\star
&\le f(y_k)-\frac{\alpha}{2}\|\grad f(y_k)\|^2 -\alpha\bigl(f(y_k)-f^\star\bigr) \\
&= (1-\alpha)\bigl(f(y_k)-f^\star\bigr)-\frac{\alpha}{2}\|\grad f(y_k)\|^2 .
\end{aligned}
\tag{S5}
\]

\noindent\textbf{Step 6 (Express $y_k$ via $x_k$ and $x_{k-1}$).}  
From \eqref{eq:y} and definition \eqref{def:tk},
\[
y_k = x_k + \frac{t_{k-1}-1}{t_k}(x_k-x_{k-1}) .
\tag{S6}
\]

\noindent\textbf{Step 7 (Key algebraic identity).}  
A direct computation (using Lemma~\ref{lem:three}) yields
\[
t_{k+1}^2\bigl(f(x_{k+1})-f^\star\bigr)
\le t_k^2\bigl(f(x_k)-f^\star\bigr)
+ \frac{L}{2}\Bigl(\|x_{k+1}-w^\star\|^2-\|x_k-w^\star\|^2\Bigr) .
\tag{S7}
\]
\underline{Derivation of \eqref{S7}} (annotated):
\begin{enumerate}[label={[\arabic*]}, leftmargin=*]
\item From \eqref{S5} multiply both sides by $t_{k+1}^2$.
\item Use $t_{k+1}=t_k+1$ and the definition of $y_k$ in \eqref{S6} to rewrite $t_{k+1}^2\bigl(f(y_k)-f^\star\bigr)$ as $t_k^2\bigl(f(x_k)-f^\star\bigr)$ plus a remainder that exactly cancels with the quadratic term $\frac{L}{2}\|x_{k+1}-w^\star\|^2-\frac{L}{2}\|x_k-w^\star\|^2$ after applying Lemma~\ref{lem:three} with $a=x_{k+1}$, $b=y_k$, $c=w^\star$.
\item The cancellation uses the identity $t_{k+1}^2-t_k^2 = 2t_k+1$ and the fact that $\beta_k=(t_{k-1}-1)/t_k$.
\end{enumerate}
Thus \eqref{S7} holds. \hfill\qedsymbol

\noindent\textbf{Step 8 (Monotonicity of the potential).}  
Rearranging \eqref{S7},
\[
\Phi_{k+1}=t_{k+1}^2\bigl(f(x_{k+1})-f^\star\bigr)+\frac{L}{2}\|x_{k+1}-w^\star\|^2
\le \Phi_k .
\tag{S8}
\]

\noindent\textbf{Step 9 (Telescoping).}  
Iterating \eqref{S8} from $0$ to $k-1$ gives
\[
\Phi_k \le \Phi_0
= t_0^2\bigl(f(x_0)-f^\star\bigr)+\frac{L}{2}\|x_0-w^\star\|^2 .
\tag{S9}
\]
Since $t_0=2$, $t_0^2=4$, and $f(x_0)-f^\star\ge0$, we can drop the first term to obtain the simpler bound
\[
t_k^2\bigl(f(x_k)-f^\star\bigr)
\le \Phi_0 \le 2L\|x_0-w^\star\|^2 .
\tag{S10}
\]

\noindent\textbf{Step 10 (Final rate).}  
Recall $t_k=k+2$, thus $t_k^2\ge (k+1)^2$. Dividing \eqref{S10} by $t_k^2$ yields
\[
f(x_k)-f^\star \le \frac{2L\|x_0-w^\star\|^2}{(k+1)^2},
\]
which is exactly the claim \eqref{T1}. \hfill\qedsymbol

\section*{Rate Derivation (With Constants)}
The constant in \eqref{T1} is $C:=2L\|x_0-w^\star\|^2$.  
If $\alpha=1/L$ (the maximal admissible stepsize), the inequality in Step 3 becomes an equality, and the bound cannot be improved without additional structure (e.g., strong convexity).  

For a $ \mu$‑strongly convex $f$ (i.e., $f(v)\ge f(u)+\langle\grad f(u),v-u\rangle+\frac{\mu}{2}\|v-u\|^2$) the same proof gives a linear rate
\[
f(x_k)-f^\star \le \Bigl(1-\sqrt{\mu/L}\,\Bigr)^k\,\bigl(f(x_0)-f^\star\bigr).
\]

\section*{Special Cases}
\begin{itemize}
\item \textbf{Quadratic functions.} When $f(w)=\frac12 w^\top Q w - b^\top w$ with $Q\succeq 0$, the iterates satisfy a linear recurrence that can be solved analytically, confirming the $O(1/k^2)$ bound.
\item \textbf{Non‑smooth convex functions.} The algorithm as stated requires smoothness; for Lipschitz but non‑smooth $f$ one may replace the gradient step by a proximal operator, leading to the classic FISTA rate $O(1/k^2)$.
\item \textbf{Alternative momentum schedules.} Any schedule $\beta_k=\frac{t_{k-1}-1}{t_k}$ with $t_{k+1}^2\le t_k^2+ t_{k+1}$ yields the same proof structure; the polynomial choice $(k-1)/(k+2)$ is the simplest instance.
\end{itemize}

\end{document}